\section{Conclusions and outlook}\label{sec:conclusion}

The Nyquist/argument-principle criterion for linear autonomous time-delay
systems was reformulated as a scalar ordinary differential equation for the
phase of the characteristic function and marched by a high-order adaptive
integrator with an exact, automatically differentiated integrand. The
outcome is a genuinely fast tool: complete stability charts of simpler
systems are computed in fractions of a second on a desktop machine --
near-real-time, fast enough for interactive exploration and for use inside
optimization loops. The reformulation tames the practically unbounded
frequency range and the near-singular integrand through the mature
step-size control of modern ODE software, and it makes the accuracy
self-validating through the integer residual of the computed winding
number, a dimensionless error measure that scales with the requested
tolerance -- empirically as $\varepsilon\approx100\,\mathrm{tol}$ -- so a
single tolerance transfers between systems with no problem-specific tuning.
It should be said plainly that the phase ODE carries no stiffness: its
right-hand side is independent of the state, so the formulation is a
quadrature in disguise, and adaptive quadrature is fully competitive for
the count alone. What the ordered march uniquely provides -- and what every
chart in this paper is built from -- is the characteristic root obtained
along the way: a smooth first-order estimate at negligible cost, sharpened
toward machine precision by an optional Taylor, Newton or combined polish
when an accurate root is wanted.

The root estimate transforms the integer-valued root-counting map into an
interpolable chart: one coloring simultaneously shows the number of unstable
roots and the spectral gap of the stable domain; the estimated root distance
serves as a scalar objective for the multi-dimensional bisection method,
which traces the true stability boundary (rather than all root-crossing
D-curves) at high equivalent resolution; the most robust parameter point
follows from the largest inscribed circle; and shifted integration lines
yield guaranteed decay-rate ($\gamma$-stability) contours. The
characteristic function itself is extracted algorithmically from a
simulation-ready right-hand side -- including nonlinear models (linearized
automatically) and singular mass matrices -- so the complete
workflow requires nothing beyond the model and the parameter ranges. The
case studies demonstrate the identical code path on retarded, neutral,
distributed-delay, differential--algebraic, transcendental,
high-dimensional and fractional-order systems, at chart-generation speeds
that allow interactive exploration and optimization.

One modelling lesson deserves to be carried away from the case studies,
because it costs nothing and outweighs any tuning of the numerics: where a
model admits a natural open-loop/feedback split, the characteristic function
should be written as the return difference $1+L(\lam)$ rather than as a raw
determinant. The zeros are the same, the added poles are the open-loop roots
and lie in the left half-plane for an open-loop stable structure, and the
leading order becomes exactly zero -- which removes both the need to estimate
it and the astronomical magnitudes that a high-order determinant produces.

The limitations delineate the outlook. First, the method is restricted to
autonomous (time-invariant) linear systems: for time-periodic coefficients
the characteristic function would have to be replaced by a Hill-type
infinite determinant, and time-domain discretization methods such as
semi-discretization~\cite{insperger2011semidiscretization,bachrathy2026multiplication}
remain the more practical choice there. Second, away from the stability
boundary the rightmost-root estimate is of first order unless refined, so
the method complements rather than replaces full spectral solvers when the
deep spectrum is needed. Third -- intrinsic to the argument principle
itself -- a peak that no sample lands on shifts the count by one while
leaving it an integer, so the residual cannot see it; the cross-check
between the count and the independently tracked root detects this at no
cost, and the optional detect-and-repair callback of
Section~\ref{sec:showcase-tolerance} mends the flagged pixels within the
march. A formally guaranteed count would still require a different device,
such as interval arithmetic. Fourth, for neutral systems beyond the
essential instability threshold the root count diverges with the truncation
frequency; the classification into stable and unstable domains nevertheless
remains correct (Appendix~\ref{app:gallery}). Fifth, the fractional-order
results are preliminary: the counting on the principal branch was validated
numerically -- including the reproduction of a published
fractional-controller design chart -- but a rigorous treatment of
branch-cut effects, in particular for shifted integration lines, is left
for future work. Finally, for large models the cost is dominated by each
evaluation of $\Dfun$ rather than by the method; where the feedback has low
rank the matrix determinant lemma already removes most of that
cost~\eqref{eq:detlemma}, and caching factorizations across frequencies,
extending the automatic extraction to neutral right-hand sides, and
coupling the interactive charts to formal optimization loops are natural
continuations.

\section*{Software and reproducibility}
The method is implemented in the open-source Julia package
\pkg{InterpolatedNyquist.jl} (MIT license,
\url{https://github.com/bachrathyd/InterpolatedNyquist.jl}), built on
\pkg{OrdinaryDiffEq.jl}~\cite{rackauckas2017differentialequations},
\pkg{ForwardDiff.jl}~\cite{revels2016forward},
\pkg{QuadGK.jl}~\cite{johnson2013quadgk},
\pkg{MDBM.jl}~\cite{bachrathy2025mdbm} and
\pkg{DelaunayTriangulation.jl}~\cite{vandenheuvel2024delaunay}. All figures and tables of this
paper can be regenerated by the scripts in the \code{paper/} directory of the
repository; the benchmark hardware and software versions are embedded
automatically in the manuscript.

\section*{Funding}
The research leading to these results was supported by the Hungarian
Scientific Research Fund (NKFIH-AG-152125).

\section*{Declaration of competing interest}
The author declares that there is no known competing financial interest or
personal relationship that could have appeared to influence the work
reported in this paper.